% ============================================================================
% 3.1 start.s
% ============================================================================
\subsection{极简的裸机程序}

这是一个完整的"在裸机上打印一句话然后永远等待"的程序。

\begin{lstlisting}[style=asm,caption=QEMU/start.s — 最简裸机程序]
.section .text
.global _start

_start:
    ldr x0, =0x09000000      ; UART0 的基地址（QEMU virt 机器的约定）
    ldr x1, =msg              ; 字符串所在的地址

print_loop:
    ldrb w2, [x1], #1        ; 读 x1 指向的字节，然后 x1++
    cbz  w2, halt             ; 读到 0（字符串结尾）就退出

wait_uart:
    ldr  w3, [x0, #0x18]     ; 读 UART 的 FLAG 寄存器 (UARTFR, 偏移 0x18)
    tbnz w3, #5, wait_uart   ; bit 5 = TXFF (发送 FIFO 满)，忙等直到不忙
    str  w2, [x0]            ; 把字节写到 UART 的数据寄存器 (UARTDR, 偏移 0x0)
    b    print_loop

halt:
    wfe                      ; 等待事件：进入低功耗状态
    b    halt                ; 万一被唤醒，继续等待

.section .rodata
msg:
    .asciz "Hello QEMU AArch64 bare metal!\n"
\end{lstlisting}

这不到 20 行做了下面几件事：

\paragraph{步骤 1：把两个关键地址取到寄存器}
\begin{itemize}
    \item \texttt{x0 = 0x0900\_0000} —— QEMU 模拟的 ARM PL011 UART0 的
        MMIO 基地址。只要往这个地址写数据，QEMU 就会把它当成字符显示到终端。
    \item \texttt{x1 = msg} —— 字符串在程序内存中的起始地址。
\end{itemize}

\paragraph{步骤 2：逐字节打印}
每一个字符的处理都要先\textbf{查询 UART 是否空闲}：

\begin{itemize}
    \item 读 \texttt{[x0, \#0x18]}（UARTFR 寄存器），检查第 5 位 (TXFF)。
    \item TXFF = Transmit FIFO Full。为 1 表示发送缓冲区已满，
        CPU 必须等待。
    \item 用 \texttt{tbnz w3, \#5, wait\_uart} 忙等这一位变 0。
    \item 一旦空闲，\texttt{str w2, [x0]} 把当前字节写出去。
\end{itemize}

PL011 是 ARM 官方定义的 UART 硬件规范，QEMU 按这个标准实现了
一个虚拟版本。所以你只要按照 PL011 的寄存器表操作，就能驱动它。

\paragraph{步骤 3：停机}
程序打印完字符串（遇到 \texttt{0} 字节即 \texttt{.asciz} 的字符串尾）后，
执行 \texttt{wfe}（Wait For Event）。这是一条 ARM 指令，
让 CPU 停止取指、进入低功耗状态，直到外部中断唤醒它。
我们没有任何中断源被启用，所以它实际上就是永久停机了。

\subsection{和 ISA 章节的差异}

与 ISA/ 目录里的用户态汇编相比，这段代码有几个不同之处：

\begin{itemize}
    \item 没有 ADRP/ADD 组合——因为我们不要求位置无关代码。
        链接器会把 \_start 放到 \texttt{0x4000\_0000}，
        所以直接用 \texttt{ldr reg, =address} 就能拿到绝对地址。
    \item 没有 \texttt{svc \#0x80}——我们直接操作硬件，没有内核帮你做事。
    \item 入口是 \texttt{\_start} 不是 \texttt{\_main}——链接脚本
        指定入口点为 \texttt{\_start}，直接在这里开始执行。
    \item \texttt{.section .rodata} 把只读数据明确放在只读段里。
\end{itemize}

整个程序唯一真正"执行"的部分就是打印——它不会从用户键盘读，
也不会做更复杂的事。后面你想扩展它的话，就是在 \texttt{halt} 之前
加更多功能（比如轮询 UART 的接收寄存器 \texttt{[x0, \#0x00]} 从用户键盘读）。
